% BHU PYQ -- Manually transcribed & error-corrected
% Paper: GLB-401: Palaeontology and Stratigraphy
% Exam: B.Sc. (Hons.) Semester IV Examination, 2022-23

\documentclass[11pt,a4paper]{article}
\usepackage[a4paper,top=2.5cm,bottom=2.5cm,left=2.8cm,right=2.8cm,headheight=14pt]{geometry}
\usepackage{amsmath,amssymb}
\usepackage[T1]{fontenc}
\usepackage[utf8]{inputenc}
\usepackage{lmodern,microtype,fancyhdr,enumitem,hyperref,lastpage,parskip}

\hypersetup{colorlinks=true,urlcolor=black,linkcolor=black,
  pdftitle={GLB-401 Palaeontology and Stratigraphy -- BHU B.Sc. Sem IV 2022-23}}
\pagestyle{fancy}\fancyhf{}
\renewcommand{\headrulewidth}{0.4pt}\renewcommand{\footrulewidth}{0.4pt}
\fancyhead[L]{\small   \textit{Banaras Hindu University}}
\fancyhead[R]{\small   \textit{GLB-401: Palaeontology \& Stratigraphy}}
\fancyfoot[C]{\small Page  \thepage\ of \pageref{LastPage}}


\newlist{parts}{enumerate}{2}
\setlist[parts,1]{label=(\roman*),leftmargin=*,itemsep=5pt,topsep=3pt}

\newcommand{\pts}[1]{\hfill   \textit{[#1]}}


\begin{document}

\begin{center}
  {\large\bfseries BANARAS HINDU UNIVERSITY}\\[2pt]
  {
\normalsize B.Sc. (Hons.) --- Semester IV Examination, 2022-23}\\[6pt]
  \rule{0.8\linewidth}{0.5pt}\\[6pt]
  {\large\bfseries Paper: GLB-401 --- Palaeontology and Stratigraphy}\\[4pt]
  {
\normalsize Time: 3 Hours \hfill Maximum Marks: 70}\\[2pt]
  {\small Ref. No.: 058430}
\end{center}
\rule{\linewidth}{0.4pt}
\smallskip



\noindent   \textit{Instructions: Answer   \textbf{five} questions in all, selecting at least   \textbf{two} each from Sections A and B, and Section-C which is compulsory. All questions carry equal marks.}
\bigskip

\bigskip


\noindent {\large\bfseries SECTION --- A}
\rule{\linewidth}{0.4pt}\smallskip




\noindent   \textbf{Question 1.} Attempt   \textbf{any two} of the following: \pts{$2  \times 7 = 14$}

\begin{parts}
  \item Discuss different conditions for fossilization.
  \item With neat sketches, describe the morphology of rhabdosome in graptolites.
  \item With neat sketches, discuss the detailed morphology of anthozoa.
\end{parts}

\medskip


\noindent   \textbf{Question 2.} With neat sketches, discuss the morphology and geological distribution of class echinoidea. \pts{14}

\medskip


\noindent   \textbf{Question 3.} Explain in detail the Geological Time Scale. \pts{14}

\medskip


\noindent   \textbf{Question 4.} Briefly explain the Cuddapah Supergroup of rocks with special emphasis on lithostratigraphic classification, lithology, and economic importance. \pts{14}

\bigskip


\noindent {\large\bfseries SECTION --- B}
\rule{\linewidth}{0.4pt}\smallskip




\noindent   \textbf{Question 5.} Write a brief note on   \textbf{any two} of the following: \pts{$2  \times 7 = 14$}

\begin{parts}
  \item Mineralization
  \item Geological distribution of trilobites
  \item Economic importance of Dharwar Supergroup
\end{parts}

\medskip


\noindent   \textbf{Question 6.} Write a brief note on   \textbf{any two} of the following: \pts{$2  \times 7 = 14$}

\begin{parts}
  \item Cephalon in trilobites with neat sketches.
  \item Principles of Stratigraphy.
  \item Physical subdivision of India with a suitable map.
\end{parts}

\medskip


\noindent   \textbf{Question 7.} Describe in detail the stratigraphy of the Aravalli craton of India with special reference to lithostratigraphic classification and economic importance. \pts{14}

\bigskip


\noindent {\large\bfseries SECTION --- C}
\rule{\linewidth}{0.4pt}\smallskip




\noindent   \textbf{Question 8.} Answer the following in 1 to 5 sentences:

\begin{parts}
  \item Difference between reworked fossil and remain fossil. \pts{2}
  \item Genal angle in trilobites. \pts{2}
  \item Index fossil. \pts{2}
  \item Calyx in Anthozoa. \pts{2}
  \item Write the name of a barite deposit located in Cuddapah Basin. \pts{1}
  \item What is gondite? \pts{1}
  \item What is the geological age of Rohtas Limestone? \pts{1}
  \item Write the name of a gold deposit located in Dharwar Craton. \pts{1}
  \item Write the names of two peninsular rivers. \pts{1}
  \item Write the name of the oldest rock of the Vindhyan Basin. \pts{1}
\end{parts}

\vfill\rule{\linewidth}{0.4pt}

\begin{center}\small
\href{https://bhu-exam.vercel.app/}{Click here to visit our website}\\[4pt]
\href{https://me-aryan.vercel.app/}{Click here to visit our contributor}\\[4pt]
\href{https://quashan-me.vercel.app/}{Click here to visit our contributor}
\end{center}
\end{document}
